% NS-54 / Astra-2. Analytic review and conditional measure-theoretic results.
% Pinned base: f0178e0; team claim: ccd722b. No arithmetic CLT is asserted.
% Append bibliography_additions.tex within the bibliography when integrating.
\subsection*{Selberg laws and the prescribed one-sided cost (NS-54)}
\label{ns54-selberg-sec}

\paragraph{Status and ranked findings.}
The proposed unconditional transfer from Selberg's central limit theorem
to the arithmetic QG input is not established. The conditional
measure-law implication below is proved, but its arithmetic hypothesis
remains an open input. The current live label is
Assumption~\ref{ass:v149-qg}, not the retired v1.48 candidate.
\begin{enumerate}
\item \textbf{MAJOR: the observable is changed.} The complete
$\beta_a=\mathcal Z_a$ is not $S(t)$ or its simple smoothing.
Its exact prime coefficients have an additional logarithmic factor;
the zero-side kernel, centering, and off-line weights must be retained.
\item \textbf{MAJOR: the cited limit law does not cover the requested
family.} The prescribed $N(a)$, observation measure, simultaneous
$a$--height limit, changing kernel, and complete remainder have not
been matched to a primary theorem. Thus neither arithmetic QG growth
nor failure of arithmetic ZLD follows from the cited CLT alone.
\item \textbf{MAJOR if full QG or a minorant obstruction is claimed:}
even a proved unweighted value law supplies no source-admissible
bounded-energy packet satisfying the all-Borel weighted domination
in \eqref{eq:v149-wlh}.
\item \textbf{MINOR, rate wording:} ``cost grows like the Gaussian
spread'' is not an asymptotic supplied by a CLT. For the actual
anchored one-sided definition, a full Gaussian limit implies the
stronger but nonquantitative divergence of cost divided by that spread
for each fixed number of levels.
\end{enumerate}
These findings concern the proposed deductions, not defects in the
existing exact identities or a proof that no future transfer can work.

\paragraph{Exact object, measure, and direction of approximation.}
Fix the entire cofinal family and its integer rule $N(a)$, and keep
\[
 L=2a,\quad x=e^{2a},\quad X_a=\pi N(a)/a,\quad
 U_a\sim\operatorname{Unif}[0,X_a],\quad
 D_a=-\inf_{\xi\in\mathbb R}\beta_a(\xi)>0.
\]
The negative pushforward is a subprobability measure:
$\widetilde m_a(B)=\Pr(\beta_a(U_a)\in B)$ for
$B\subset[-D_a,0]$. It is not normalized conditionally on negativity.
The admissible grids contain $-D_a$, have at most $K$ levels in
$[-D_a,0]$, and round \emph{down}:
\[
 p_{\mathcal L}(t)=\max\{\ell\in\mathcal L:\ell\le t\},\qquad
 \mathrm{QC}_K=\inf_{\mathcal L}\int(t-p_{\mathcal L}(t))\,d\widetilde m_a(t).
\]
The depth is global, not the observed head depth or the depth of a
smoothed zero prefix. All zeros, endpoint and trivial terms, and the
symmetric-height Perron convention in
Proposition~\ref{prop:v146-beta-explicit} remain in the field.

\paragraph{Exact identity: a differentiated finite logarithmic polynomial.}
Put $A(\xi)=\Re\psi(5/4+i\xi/2)-\log\pi$ and define
\[
 V_x(\xi)=-\frac1\pi\sum_{2\le n<x}
       \frac{\Lambda(n)}{\sqrt n\log n}\sin(\xi\log n),\qquad
 I_x(\xi)=\Re\frac{x^{1/2+i\xi}-1}{1/2+i\xi}.
\]
The exact prime expression gives, by finite differentiation,
\begin{equation}\label{ns54-selberg-eq-derivative}
             \beta_a(\xi)=A(\xi)+2\pi V_x'(\xi)+2I_x(\xi).
\end{equation}
No approximation of $S$ by $V_x$ is asserted. Even such an
approximation without derivative control would not suffice:
$n^{-1}\sin(n^2t)\to0$ uniformly whereas its derivative does not.
Convergence of finitely many values likewise gives no continuity
theorem for a growing differential or convolution operator.

\begin{proposition}[Fixed-window long-height calibration]
\label{ns54-selberg-prop-fixed}
For fixed $a$, let $B_a=\beta_a-A$ and average uniformly on $[T,2T]$.
Then
\begin{equation}\label{ns54-selberg-eq-fixed}
 \mathbb E_T B_a=O_a(T^{-1}),\qquad
 \operatorname{Var}_T B_a
   =2\sum_{2\le n<e^{2a}}\frac{\Lambda(n)^2}{n}+O_a(T^{-1}).
\end{equation}
In particular the variance in this limit is not a universal
$\log\log T$ variance. Also $\beta_a(\xi)=\log\xi+O_a(1)$ as
$\xi\to+\infty$.
\end{proposition}
\begin{proof}
Equation~\eqref{ns54-selberg-eq-derivative} says
$B_a=-2\sum_{n<x}\Lambda(n)n^{-1/2}\cos(\xi\log n)+2I_x$.
For fixed $x$, $I_x=O_a(T^{-1})$ on $[T,2T]$.
The average of each cosine is $O_a(T^{-1})$; the distinct
frequencies $\log n$ have vanishing cross averages with this same
fixed-$a$ error. Each square has average $1/2+O_a(T^{-1})$.
This proves both moment assertions. The finite sum is bounded and
$A(\xi)=\log(\xi/(2\pi))+o(1)$, proving the last assertion.
No uniformity in growing $a$ is claimed.
\end{proof}

\paragraph{Zero-side interpretation, explicitly conditional where necessary.}
Only under RH may every zero be replaced by its ordinate in the
real sinc expression. Writing $k_L(u)=\sin(Lu)/u$, with $k_L(0)=L$,
that expression is
\[
 \beta_a(\xi)=2\lim_{H\to\infty}
              \sum_{|\gamma|<H}k_L(\gamma-\xi)+C_a(\xi).
\]
It includes the density contribution; it is not the negative of a
centered density fluctuation. On a finite positive-height interval
away from endpoint zeros, $N(t)=\vartheta(t)/\pi+1+S(t)$ gives
\[
 \int_A^B k_L(t-\xi)\,dN(t)
 =\int_A^B k_L(t-\xi)\vartheta'(t)/\pi\,dt
  +[k_L(t-\xi)S(t)]_A^B
  -\int_A^B k_L'(t-\xi)S(t)\,dt.
\]
Thus the centered part involves a derivative kernel, with boundary
and limiting terms. Unconditionally the complete formula instead
contains $x^{\rho-1/2}$ and $\rho-1/2+i\xi$; discarding real parts
of zeros is not licensed by a theorem about their ordinates.

\paragraph{What the primary CLTs actually say.}
Radziwi\l\l--Soundararajan's Theorem~1 is unconditional for
$\log|\zeta(1/2+it)|/\sqrt{\tfrac12\log\log T}$ with
$t$ uniform on $[T,2T]$ and fixed thresholds
\cite{ns54-RadziwillSoundararajan}.
Bourgade's Theorem~1.1 supplies joint laws for a fixed number of
shifted logarithms with prescribed separation limits, not a
growing derivative operator \cite{ns54-Bourgade}.
Maples--Rodgers' Theorem~1.2 is genuinely unconditional: it treats
centered ordinate counts with a fixed compactly supported
bounded-variation test, $n(T)\to\infty$, $n(T)=o(\log T)$, and
diverging Fourier variance \cite{ns54-MaplesRodgers}.
The sinc profile is not that test: it has noncompact support,
infinite total variation, and finite $H^{1/2}$ Fourier energy.
Bourgade--Kuan's Theorem~1 assumes RH and additional decay and
regularity \cite{ns54-BourgadeKuan}; its stated decay also excludes
the unmodified sinc profile. These observations delimit those
theorems, not all possible extensions of their methods.

\begin{proposition}[Quantitative sufficient transfer to fixed-level growth]
\label{ns54-selberg-prop-transfer}
Let $b_a\to\infty$, $c\in\mathbb R$, and suppose the distribution
function of $Y_a=\beta_a(U_a)/b_a$ has distance at most $\Delta_a$
from $z\mapsto\Phi(z-c)$, where $\Delta_a\to0$ and $\Phi$ is
the standard Gaussian distribution function. For fixed $K$, put
\[
 J_j=(-3j-1,-3j),\quad 1\le j\le K+1,\qquad
 m_K=\min_j\{\Phi(-3j-c)-\Phi(-3j-1-c)\}>0.
\]
Eventually $\Delta_a\le m_K/4$ and
\begin{equation}\label{ns54-selberg-eq-transfer}
                  \mathrm{QC}_K(\widetilde m_a)\ge b_am_K/2.
\end{equation}
An explicit sufficient surrogate hypothesis is: for real $F_a$,
centers $\mu_a$, and nonnegative $\delta_a,r_a,\epsilon_a$,
\[
 \sup_z\left|\Pr\left(\frac{F_a(U_a)-\mu_a}{b_a}\le z\right)
                    -\Phi(z)\right|\le\delta_a,\qquad
 \Pr(|F_a(U_a)-\beta_a(U_a)|>r_ab_a)\le\epsilon_a.
\]
One can then take
\begin{equation}\label{ns54-selberg-eq-error}
 \Delta_a=\delta_a+\epsilon_a+
             \frac{r_a+|\mu_a/b_a-c|}{\sqrt{2\pi}}.
\end{equation}
\end{proposition}
\begin{proof}
The Gaussian density is at most $1/\sqrt{2\pi}$. Sandwiching the
two distribution functions outside the exceptional event proves
\eqref{ns54-selberg-eq-error}. Each $J_j$ has $Y_a$-mass at least
$m_K-2\Delta_a\ge m_K/2$; continuity of the actual field's value
distribution removes endpoint atoms. Their intersections with its
support retain that mass and the pairwise distance is at least $2$.
For any admissible grid the set on which downward error is less
than $b_a$ lies in at most $K$ intervals of length $b_a$.
Each meets at most one of the $K+1$ separated bands $b_aJ_j$.
On a missed band the loss is at least $b_a$, proving the bound.
\end{proof}

\begin{proposition}[Full unbounded-left laws force a stronger one-sided cost]
\label{ns54-selberg-prop-unbounded}
Suppose, on the prescribed entire family,
$\beta_a(U_a)/b_a\Rightarrow Y$ with $b_a>0$, and the probability
law of $Y$ has unbounded support to the left. Then
\begin{equation}\label{ns54-selberg-eq-unbounded}
 \frac{D_a}{b_a}\longrightarrow\infty,\qquad
 \frac{\mathrm{QC}_K(\widetilde m_a)}{b_a}\longrightarrow\infty
                   \quad\hbox{for every fixed }K\ge1.
\end{equation}
If also $b_a\to\infty$, the QG growth clause follows. On the same
family ZLD fails for every fixed $\theta>0$.
\end{proposition}
\begin{proof}
For any $M>0$ the limiting law gives positive mass to an open
interval below $-M$. Weak convergence therefore implies
$D_a/b_a>M$ eventually. This proves the first limit.
If the second failed along a subsequence, choose grids whose
normalized costs are bounded, within $1$ of the infima.
Order the normalized levels and pad by repetitions to $K$ entries.
By compactness of the extended interval $[-\infty,0]$, pass to a
subsequence on which every entry converges. At least the entry
$-D_a/b_a$ tends to $-\infty$.

If all entries tend to $-\infty$, choose a bounded negative open
interval with positive limiting mass. All downward losses there
tend uniformly to infinity. Otherwise let $\ell_*$ be the smallest
finite limiting level and choose such an interval strictly below
$\ell_*$. Every finite-limit level is eventually above the interval;
every other level tends to $-\infty$. Hence again the loss there
tends uniformly to infinity. Weak convergence gives a positive
lower bound for its mass, contradicting bounded normalized cost.
All omitted losses are nonnegative, so no moment-convergence
assumption is needed.

Finally weak convergence makes $\beta_a(U_a)/b_a$ tight. For
$0<r_1<r_2\le\theta$,
\[
 \Pr(-r_2D_a<\beta_a(U_a)<-r_1D_a)
 \le\Pr(\beta_a(U_a)/b_a<-r_1D_a/b_a)\longrightarrow0.
\]
This contradicts ZLD's fixed positive lower bound
$\kappa(r_2-r_1)$. The conclusion is conditional on the full
head-sampling law; a Gaussian law on just one subinterval does
not control the remaining head mass.
\end{proof}

\paragraph{Why centering and negative mass cannot be omitted.}
A CLT only for $(\beta_a-\mu_a)/b_a$ leaves $\mu_a/b_a$ unknown.
As a measure-theoretic countermodel, take $Z_n$ standard Gaussian
conditioned on $|Z_n|\le n$ and $Y_n=b_n(2n+Z_n)$ with
$b_n\to\infty$. Its centered normalized law is Gaussian, but it
has no negative mass and hence zero negative cost. A negative
global minimum may be placed outside the observation interval
without changing this example. This is not an arithmetic
counterexample. Similarly, a limit conditional on negative values
has to restore their fraction $\varphi_-(a)$: separated bands then
give a cost of at least a constant times $b_a\varphi_-(a)$, and
that product must diverge. A positive limiting negative fraction
follows from a full Gaussian law with finite normalized centering;
it is not a consequence of Gaussian fluctuation language alone.

\paragraph{Exact open transfer lemma and scope.}
To use Selberg-type methods here, prove the surrogate hypotheses of
Proposition~\ref{ns54-selberg-prop-transfer}, or another sufficient
band-law hypothesis, for the \emph{actual} prescribed $N(a)$.
Specify height $T(a)$ and the relation of the kernel width $1/(2a)$
to mean zero spacing $2\pi/\log T(a)$. Matching the scale
$L=2a$ to Maples--Rodgers' $\log T/(2\pi n(T))$ requires
$n(T)=\log T/(2\pi L)$. Its required divergence fails when
$\log T\asymp L$; the theorem cannot then be used by silently
choosing a different head rule. A law on $[T,2T]\subset[0,X_a]$
contributes only the factor $T/X_a$ to unweighted band masses.
Thus a vanishing sampling fraction may destroy the growth transfer;
NS-40's exact cutoff dilution still applies.

An ordinate-only surrogate must control the omitted off-line
weights and shifted denominators. A finite-zero surrogate must
control $-2\Re R_{x,H}(1/2-i\xi)$ with $H>X_a$, as well as
smoothing/desmoothing errors; that inequality on $H$ is not an
error bound. The endpoint/trivial correction alone is uniformly
bounded by $e_x+4x^{-5/2}/(5(1-x^{-2}))\to0$, which is adequate
relative to $b_a\to\infty$ but is not the Perron remainder.
Errors $o(b_a)$ outside a fraction $o(1)$ suffice for the displayed
Gaussian transfer if the surrogate law and centering are proved.

None of these arithmetic distribution and error hypotheses is
established by this review. No particular cofinal $N(a)$ is supplied
by the bare definition of QG. The packet requirement remains
separate even after the growth clause. There is no new numerical
window, signed complementary floor, G2, or RH conclusion.
